\documentclass[10pt,a4paper]{article}
\usepackage{amsmath,amssymb,amsthm,graphicx}
\usepackage{amsmath,amssymb,amsthm}
\newcommand{\supp}{\operatorname{supp}}
\newcommand{\Newt}{\operatorname{Newt}}
\newcommand{\conv}{\operatorname{conv}}
\newcommand{\vertset}{\operatorname{vert}}
\newenvironment{sourceresult}[1]{\par\medskip\noindent\textbf{#1.}\itshape}{\par\medskip}
\title{Shadow complexity lower bounds monotone circuit size for transparent polynomials}
\author{Pavel Hrubeš, Amir Yehudayoff}
\date{doi:10.4230/LIPIcs.CCC.2021.9; CC BY 4.0}
\begin{document}
\maketitle
\noindent\textit{Editorial note.} Faithful editable transcription of official CC BY publication; author-hosted/ECCC pages searched for native TeX without locating a public source archive. Author numbering and original page locators retained. Standard planar Minkowski background retained with the author-provided outline rather than expanded proof. No mathematical correction or added author argument. Actual journal edition is an identity alias only; its additional content is not used. Production and validation only on Yupeng.\par\medskip
\noindent\textit{Allowed standard prerequisites (editorial).}\begin{itemize}
\item Convex hulls, finite convexly independent sets and arithmetic circuits. Standard background; original shadow, transparency, h.p.-circuit and h.p.-projection definitions retained. Folklore planar Minkowski-sum lemma retained with author outline and original cited source.
\end{itemize}\medskip
\section*{Original definitions and computational models (pp.9:1--2,9:11--14)}
A polytope is the convex hull of a finite set of points in Euclidean space. Equivalently, it is a compact set that is defined by finitely many linear inequalities.
A shadow of a polytope $P\subseteq\mathbb R^n$ is a set of the form $L(P)$, where $L:\mathbb R^n\to\mathbb R^2$ is a linear map. The shadows of $P$ are two-dimensional polygons, and hence typically much simpler than $P$. The shadow complexity of $P$ is
\[
\sigma(P)=\max_L |\vertset(L(P))|,
\]
where $L$ is a linear map and $\vertset(Q)$ is the vertex set of the polytope $Q$.
Let $\mathbb F$ be a field. For a list of variables $x=(x_1,\ldots,x_n)$ and $\alpha\in\mathbb N^n$, let $x^\alpha$ be the monomial $\prod_{i=1}^n x_i^{\alpha_i}$. A polynomial $f\in\mathbb F[x_1,\ldots,x_n]$ is a formal sum of the form $\sum_{\alpha\in\mathbb N^n}a_\alpha x^\alpha$ where $\supp(f):=\{\alpha\in\mathbb N^n:a_\alpha\ne0\}$ is finite. The Newton polytope of $f$ is
\[\Newt(f):=\conv(\supp(f)),\]
where $\conv(\cdot)$ denotes the convex hull.
A polynomial $f$ is transparent if $|\supp(f)|=\sigma(\Newt(f))$. In other words, there is a linear map $L:\mathbb R^n\to\mathbb R^2$ which maps $\supp(f)$ to distinct convexly independent points in $\mathbb R^2$.

As the standard model of computation of polynomials we take the arithmetic circuit model. It is a (finite) directed acyclic graph whose every node has in-degree zero or two. Nodes of in-degree zero (input nodes) are labelled with a constant or a variable. Nodes of in-degree two are labelled with operations $+$ or $\times$. Every node of the circuit computes a polynomial in $\mathbb F$ in the natural way. As the size of the circuit, we take the number of nodes. A circuit is called a formula if its underlying graph is a tree. For more background and motivation, see [44] and references within.
Our focus is mainly on monotone computation. A polynomial over $\mathbb R$ is monotone if all of its coefficients are non-negative. Similarly, a monotone arithmetic circuit can use only non-negative constants.

In a Laurent polynomial, variables are allowed to have negative integer exponents. The notions of $\supp(f)$ and Newton polytope of $f$ are defined in the obvious manner. A Laurent polynomial over $\mathbb R$ is monotone, if all of its coefficients are non-negative.
A high powered circuit (h.p.-circuit for short) is an arithmetic circuit in which every input node is labelled by a term $ax_1^{k_1}\cdots x_n^{k_n}$ with $a\in\mathbb F$ and $k_1,\ldots,k_n\in\mathbb Z$. The size of the circuit is the number of its nodes.
In other words, we have given the circuit a power to compute every term $ax^\alpha$ at a unit cost. This is especially important in the case of h.p.-formula. An arithmetic formula of size $s$ can compute a polynomial of degree at most $s$, whereas there is no such restriction in an h.p.-formula. Furthermore, we have allowed the variables to have negative exponents and hence an h.p.-circuit computes a Laurent polynomial instead of a polynomial. But this is only a cosmetic detail.

A high power projection (h.p.-projection for short) is a homomorphism
\[\pi:\mathbb F[x_1,\ldots,x_n]\to\mathbb F[y_1,\ldots,y_m,y_1^{-1},\ldots,y_m^{-1}]\]
such that $\pi(x_i)=a_i y^{\alpha_i}$ for every $x_i$, where $a_i\in\mathbb F$ and $\alpha_i\in\mathbb Z^m$, and for every $f\in\mathbb F[x_1,\ldots,x_n]$,
\[\pi(f(x_1,\ldots,x_n))=f(\pi(x_1),\ldots,\pi(x_n)).\]
The constants $a_i$ are called the coefficients of $\pi$ and $\alpha_i$ its exponents. If $\mathbb F=\mathbb R$ and $\pi$ has non-negative coefficients, then $\pi$ is called monotone.
An h.p.-projection $\pi$ induces a linear map $L_\pi:\mathbb R^n\to\mathbb R^m$ by setting $L_\pi(e_i)=\alpha_i$ and extending it linearly to $\mathbb R^n$. It follows that $\supp(\pi(f))\subseteq L_\pi(\supp(f))$. The inclusion may be strict, as some monomials of $f$ can cancel out in the projection. If no cancellations occur, we indeed have $\Newt(\pi(f))=L_\pi(\Newt(f))$. This is satisfied, e.g., if $f$ is monotone and the coefficients of $\pi$ are positive, or if the coefficients are algebraically independent.

\begin{sourceresult}{Theorem 2 (p.9:2)}
If $f$ is transparent then every monotone circuit computing $f$ has size at least $\Omega(\sigma(\Newt(f)))$.
\end{sourceresult}

\section*{Original bearing geometric statements (pp.9:6,9:12,9:14--15)}
Given $A,B\subseteq\mathbb R^n$, their Minkowski sum is defined as
\[A+B:=\{a+b:a\in A,b\in B\}.\]
If $P$ and $Q$ are polytopes then $P+Q$ is also a polytope. In two-dimensions, Minkowski sum has nice properties. Let $P$ be a polytope in $\mathbb R^2$ with vertices $v_1,\ldots,v_k$ where $k>1$. We can assume they are ordered so that $P$ lies in the left closed half plane determined by the line going from $v_i$ to $v_{i+1}$ for $i<k$, and similarly for $v_k$ and $v_1$. Let $E(P)$ be the collection of unit vectors in the direction of these $k$ edges. That is, vectors of the form $(v_{i+1}-v_i)/\|v_{i+1}-v_i\|$ for $i<k$, and $(v_1-v_k)/\|v_1-v_k\|$. If $|\vertset(P)|\le1$ then $E(P):=\emptyset$.
\begin{sourceresult}{Lemma 8}
Let $P_1,\ldots,P_r$ be non-empty polytopes in $\mathbb R^2$. Then $E(P_1+\cdots+P_r)=\bigcup_{i=1}^r E(P_i)$. Consequently, $|\vertset(P_1+\cdots+P_r)|\le\sum_{i=1}^r|\vertset(P_i)|$. The latter holds for empty $P_i$'s as well.
\end{sourceresult}
The lemma is folklore. It can be inferred from Chapter 13.3 in [11], and we give only an outline of proof.
\begin{proof}[Proof sketch of Lemma 8]
Given a non-empty polytope $P$ and $w\in\mathbb R^2$, let $P^w:=\{x\in P:\langle x,w\rangle=\max_{z\in P}\langle z,w\rangle\}$ be the set of extreme of points of $P$ in the direction $w$. It is either a vertex or an edge of $P$. For a pair of polytopes we have $(P_1+P_2)^w=P_1^w+P_2^w$. Every edge of $P_1$ yields an edge of $P_1+P_2$ with the same direction. Conversely, every edge of $P_1+P_2$ comes from one of $P_1$ or $P_2$.
\end{proof}

\begin{sourceresult}{Lemma 28}
Let $f$ be a polynomial over an infinite field $\mathbb F$. Assume that $\sigma(\Newt(f))=k$. Then there exists a bivariate Laurent polynomial $g\in\mathbb F(y_1,y_2,y_1^{-1},y_2^{-1})$ which is an h.p.-projection of $f$ so that $\Newt(g)$ has $k$ vertices. Moreover, if $f$ is a homogeneous polynomial then $g$ is a polynomial. If $\mathbb F=\mathbb R$, then the coefficients in the projection can be assumed positive.
\end{sourceresult}
\begin{proof}
Let $L(z)=Az$ with $A\in\mathbb R^{2\times n}$ be a linear map so that
\[|\vertset(L(\Newt(f)))|=k.\]
We can assume that the entries of $A$ are rational, because a small perturbation of $A$ cannot decrease $|\vertset(L(\Newt(f)))|$. Now, we can assume that the entries of $A$ are integers, because we can multiply $A$ by a suitable integer.
Moreover, when $f$ is homogeneous of degree $d$, increasing all entries of $A$ by $\lambda$ corresponds to shifting $L(\Newt(f))$ by $(\lambda d,\lambda d)$, which does not change the number of vertices. Hence, in this case, $A$ can be taken with non-negative integer entries.
Let us now consider a h.p.-projection $\pi$ with $\pi(x_i)=a_i y^{L(e_i)}$. It follows that $\supp(\pi(f))\subseteq L(\supp(f))$. Now, we claim that we can choose the coefficients $a_i$ so that equality holds. This can be seen as follows. Given $\alpha\in\supp(f)$, the coefficient of $y^{L(\alpha)}$ in $\pi(f)$ is a non-zero polynomial $h_\alpha$ in the coefficients of $\pi$. Hence, if $\mathbb F$ is infinite, there exist non-zero coefficients for which $h_\alpha$ is non zero for every $\alpha\in\supp(f)$. If $\mathbb F=\mathbb R$, they can be taken positive.
\end{proof}

\begin{sourceresult}{Lemma 37}
Let $A,B\subseteq\mathbb R^2$ be non-empty sets such that $A+B$ is a convexly independent set with $|A|\ge|B|$. Then either $|A|\le2$ or $|B|\le1$.
\end{sourceresult}
\begin{proof}
For the sake of contradiction, assume that $A+B$ is convexly independent, $|A|\ge3$ and $|B|\ge2$. By Lemma 8, the convex hull of $A+B$ has at most $|A|+|B|$ vertices. By the size assumption, there exist $a_1\ne a_2\in A$ and $b_1\ne b_2\in B$ with $a_1+b_1=a_2+b_2$. The point $a_1+b_1$ is the average of $a_1+b_2$ and $a_2+b_1$ and it is distinct from them, a contradiction.
\end{proof}

\section*{Entire original circuit argument (Section 5.5, pp.9:15--16)}
\begin{sourceresult}{Theorem 38}
Let $f$ be a monotone bivariate Laurent polynomial such that $\supp(f)$ is convexly independent and $|\supp(f)|=k$. Then $f$ requires monotone h.p.-circuit with $k/4$ gates.
\end{sourceresult}
Theorem 38 implies Theorem 2 via Lemma 28.
\begin{proof}
The lower bound is proved using the following ``progress'' measure. Given $A\subseteq\mathbb R^2$ and $\epsilon\in\{0,1\}$, let $A^\epsilon:=A$ if $\epsilon=1$ and $A^\epsilon:=\emptyset$ if $\epsilon=0$. Given $v\in\mathbb R^2$, let $v+A:=\{v\}+A$. Let $\mathcal A$ be a finite set of finite subsets of $\mathbb R^2$. For functions $\epsilon:\mathcal A\to\{0,1\}$ and $v:\mathcal A\to\mathbb R^2$, let
\[\mathcal A_{\epsilon,v}=\bigcup_{A\in\mathcal A}(v(A)+A)^{\epsilon(A)}.\]
Let
\[\mu(\mathcal A)=\max_{\epsilon,v}\{|\mathcal A_{\epsilon,v}|:\mathcal A_{\epsilon,v}\text{ is convexly independent}\}.\]
\par\medskip\noindent\textbf{Claim.} Let $\mathcal A'=\mathcal A\cup\{B\}$ and $A_1,A_2\in\mathcal A$. Then
\begin{align}
\mu(\mathcal A')&\le\mu(\mathcal A)+|B|,\tag{3}\\
\mu(\mathcal A')&\le\mu(\mathcal A)+2,&&\text{if }B=u+A_1\text{ for some }u\in\mathbb R^2,\tag{4}\\
\mu(\mathcal A')&\le\mu(\mathcal A)+4,&&\text{if }B=A_1\cup A_2,\tag{5}\\
\mu(\mathcal A')&\le\mu(\mathcal A)+4,&&\text{if }B=A_1+A_2.\tag{6}
\end{align}
\begin{proof}[Proof of Claim]
Inequality (3) is straightforward.
To prove (4), suppose that $\epsilon,v$ are such that $\mathcal A'_{\epsilon,v}$ is convexly independent. Suppose $\epsilon(A_1)=\epsilon(B)=1$ and $v(A_1)+A_1\ne v(B)+B$; otherwise we have $|\mathcal A'_{\epsilon,v}|\le\mu(\mathcal A)$. Then $(v(A_1)+A_1)\cup(v(B)+B)=\{v(A_1),v(B)+u\}+A_1$ is convexly independent. Since $|\{v(A_1),v(B)+u\}|=2$, by Lemma 37, $A_1$ has size at most 2. This means $\mu(\mathcal A')\le\mu(\mathcal A)+2$ by (3).
For (5), observe that $\mu(\mathcal A')\le\mu(\mathcal A\cup\{u_1+A_1,u_2+A_2\})$ whenever $u_1,u_2\ne0$ are distinct and apply (4) twice.
Finally, we prove (6). If $B=A_1+A_2$ is not convexly dependent, it contributes nothing to $\mu$. Assume that $B$ is convexly independent and $|A_1|\ge|A_2|>0$. By Lemma 37, either $|A_1+A_2|\le4$ or $|A_2|=1$. In the former case, $\mu(\mathcal A')\le\mu(\mathcal A)+4$ by (3). In the latter, $A_2=\{u\}$ for some $u$ and $\mathcal A'=\mathcal A\cup\{u+A_1\}$ and we can apply (4).
\end{proof}
Let us call a h.p.-circuit transparent, if every gate in the circuit computes a polynomial with convexly independent support. Given a circuit $\Psi$ and a node $u$, let $\supp(u)$ be the support of the Laurent polynomial computed by $u$. Let $\mathcal A_\Psi$ be the set $\{\supp(u):u\in\Psi\}$.
Using the Claim, we can show that whenever a transparent and monotone $\Psi$ has $s$ gates then $\mu(\mathcal A_\Psi)\le4s$. The proof is by induction. The induction base $s=1$ trivially holds. It remains to verify the induction step. Let $u$ be an output gate of $\Psi$. If $u$ is also an input gate, apply (3).
If $u=u_1+u_2$ then $\supp(u)=\supp(u_1)\cup\supp(u_2)$ and (5) completes the proof.
If $u=u_1\times u_2$ then $\supp(u)=\supp(u_1)+\supp(u_2)$ and (6) completes the proof.
Finally, consider a monotone circuit $\Psi$ for $f$ of minimal size $s$. No gate in the circuit computes the zero polynomial (unless $f$ itself the zero polynomial). The circuit is transparent because a monotone computation does not cancel monomials unless multiplying by zero, and because $+,\times$ can not ``undo'' convex dependence. This means that $\mu(\mathcal A_\Psi)\le4s$. On the other hand, since $\supp(f)$ consists of $k$ convexly independent points, we have $\mu(\mathcal A_\Psi)\ge|\supp(f)|=k$.
\end{proof}

\begin{thebibliography}{44}
\bibitem[11]{source11} M. de Berg, M. van Kreveld, M. Overmars, and O. Schwarzkopf. \emph{Computational Geometry: Algorithms and Applications}. Springer, 2 edition, 2000.
\bibitem[44]{source44} A. Shpilka and A. Yehudayoff. Arithmetic circuits: A survey of recent results and open questions. \emph{Foundations and Trends in Theoretical Computer Science}, 5(3-4), 2010.
\end{thebibliography}
\end{document}
